\documentclass[aspectratio=169]{beamer}

% ============================================
% THEME AND COLOR SCHEME
% ============================================

\usetheme{Madrid}
\usecolortheme{default}

% Alternative themes: AnnArbor, Berlin, Copenhagen, Madrid, Singapore, Warsaw
% Alternative color themes: beaver, crane, dolphin, orchid, rose, seagull

% ============================================
% PACKAGE IMPORTS
% ============================================

% Encoding
\usepackage[utf8]{inputenc}
\usepackage[T1]{fontenc}

% Graphics
\usepackage{graphicx}
\graphicspath{{images/}}

% Tables
\usepackage{booktabs}

% Math
\usepackage{amsmath, amssymb}

% Code listings
\usepackage{listings}
\lstset{
    basicstyle=\ttfamily\small,
    keywordstyle=\color{blue},
    commentstyle=\color{green!50!black},
    stringstyle=\color{red},
    showstringspaces=false,
    breaklines=true,
    frame=single,
    numbers=left,
    numberstyle=\tiny\color{gray}
}

% Hyperlinks
\usepackage{hyperref}
\hypersetup{
    colorlinks=true,
    linkcolor=blue,
    urlcolor=cyan
}

% ============================================
% CUSTOM COMMANDS
% ============================================

% Custom color definitions
\definecolor{myblue}{RGB}{0,102,204}
\definecolor{mygreen}{RGB}{0,153,51}
\definecolor{myred}{RGB}{204,0,0}

% ============================================
% PRESENTATION INFORMATION
% ============================================

\title[Short Title]{Your Presentation Title}
\subtitle{Optional Subtitle}

\author[Short Name]{
    Your Name\inst{1} \and 
    Co-Author Name\inst{2}
}

\institute[Short Institute]{
    \inst{1}%
    Department Name\\
    University Name
    \and
    \inst{2}%
    Another Department\\
    Another University
}

\date[Conference 2024]{
    Conference Name 2024 \\
    \today
}

% Logo (uncomment and provide logo file)
% \logo{\includegraphics[height=1cm]{logo.png}}

% ============================================
% PRESENTATION CONTENT
% ============================================

\begin{document}

% ============================================
% TITLE SLIDE
% ============================================

\begin{frame}
    \titlepage
\end{frame}

% ============================================
% TABLE OF CONTENTS
% ============================================

\begin{frame}{Outline}
    \tableofcontents
\end{frame}

% ============================================
% SECTION 1: INTRODUCTION
% ============================================

\section{Introduction}

\begin{frame}{Introduction}
    \begin{block}{Problem Statement}
        Describe the problem or topic you are addressing.
    \end{block}
    
    \vspace{0.3cm}
    
    \begin{itemize}
        \item Key challenge 1
        \item Key challenge 2
        \item Key challenge 3
    \end{itemize}
\end{frame}

\begin{frame}{Motivation}
    \begin{columns}[T]
        \begin{column}{0.5\textwidth}
            \textbf{Why is this important?}
            \begin{itemize}
                \item Reason 1
                \item Reason 2
                \item Reason 3
            \end{itemize}
        \end{column}
        
        \begin{column}{0.5\textwidth}
            \textbf{Current limitations:}
            \begin{itemize}
                \item Limitation 1
                \item Limitation 2
                \item Limitation 3
            \end{itemize}
        \end{column}
    \end{columns}
\end{frame}

\begin{frame}{Research Questions}
    \begin{enumerate}
        \item \textbf{RQ1:} First research question?
        \item \textbf{RQ2:} Second research question?
        \item \textbf{RQ3:} Third research question?
    \end{enumerate}
\end{frame}

% ============================================
% SECTION 2: BACKGROUND
% ============================================

\section{Background}

\begin{frame}{Background and Related Work}
    \begin{block}{Key Concepts}
        \begin{itemize}
            \item Concept 1: Brief explanation
            \item Concept 2: Brief explanation
            \item Concept 3: Brief explanation
        \end{itemize}
    \end{block}
    
    \vspace{0.3cm}
    
    \begin{alertblock}{Important Note}
        Highlight important information here.
    \end{alertblock}
\end{frame}

\begin{frame}{Previous Work}
    \begin{table}
        \centering
        \begin{tabular}{lcc}
            \toprule
            \textbf{Approach} & \textbf{Pros} & \textbf{Cons} \\
            \midrule
            Method A & Fast & Limited \\
            Method B & Accurate & Slow \\
            Method C & Simple & Basic \\
            \bottomrule
        \end{tabular}
        \caption{Comparison of existing methods}
    \end{table}
\end{frame}

% ============================================
% SECTION 3: METHODOLOGY
% ============================================

\section{Methodology}

\begin{frame}{Proposed Approach}
    \begin{block}{Overview}
        Brief overview of your approach.
    \end{block}
    
    \vspace{0.3cm}
    
    \textbf{Key steps:}
    \begin{enumerate}
        \item Step 1: Description
        \item Step 2: Description
        \item Step 3: Description
        \item Step 4: Description
    \end{enumerate}
\end{frame}

\begin{frame}{Architecture/Framework}
    \begin{center}
        % Uncomment when you have a figure
        % \includegraphics[width=0.8\textwidth]{architecture.png}
        
        \vspace{1cm}
        [Place your architecture diagram here]
        \vspace{1cm}
    \end{center}
    
    \begin{itemize}
        \item Component A: Purpose
        \item Component B: Purpose
        \item Component C: Purpose
    \end{itemize}
\end{frame}

\begin{frame}[fragile]{Algorithm/Code Example}
    \begin{lstlisting}[language=Python]
def example_function(data):
    """
    Example function demonstrating code
    """
    result = []
    for item in data:
        if item > threshold:
            result.append(process(item))
    return result
    \end{lstlisting}
\end{frame}

\begin{frame}{Mathematical Formulation}
    The objective function is defined as:
    
    \begin{equation}
        \min_{x} f(x) = \sum_{i=1}^{n} (y_i - \hat{y}_i)^2
    \end{equation}
    
    Subject to constraints:
    \begin{align}
        g(x) &\leq 0 \\
        h(x) &= 0
    \end{align}
    
    where $x$ represents the decision variables.
\end{frame}

% ============================================
% SECTION 4: RESULTS
% ============================================

\section{Results}

\begin{frame}{Experimental Setup}
    \begin{columns}[T]
        \begin{column}{0.5\textwidth}
            \textbf{Datasets:}
            \begin{itemize}
                \item Dataset A (1000 samples)
                \item Dataset B (5000 samples)
                \item Dataset C (10000 samples)
            \end{itemize}
        \end{column}
        
        \begin{column}{0.5\textwidth}
            \textbf{Evaluation Metrics:}
            \begin{itemize}
                \item Accuracy
                \item Precision
                \item Recall
                \item F1-Score
            \end{itemize}
        \end{column}
    \end{columns}
\end{frame}

\begin{frame}{Performance Results}
    \begin{table}
        \centering
        \begin{tabular}{lccc}
            \toprule
            \textbf{Method} & \textbf{Acc.} & \textbf{Prec.} & \textbf{Rec.} \\
            \midrule
            Baseline & 85.3 & 82.1 & 80.5 \\
            Method A & 88.7 & 85.4 & 84.2 \\
            Method B & 90.2 & 87.9 & 86.8 \\
            \textbf{Ours} & \textbf{93.5} & \textbf{91.2} & \textbf{90.1} \\
            \bottomrule
        \end{tabular}
        \caption{Performance comparison}
    \end{table}
\end{frame}

\begin{frame}{Visual Results}
    \begin{center}
        % Uncomment when you have results
        % \includegraphics[width=0.9\textwidth]{results.png}
        
        \vspace{2cm}
        [Place your results visualization here]
        \vspace{2cm}
    \end{center}
\end{frame}

\begin{frame}{Key Findings}
    \begin{itemize}
        \item<1-> \textbf{Finding 1:} Our method outperforms baselines by 3.3\%
        \item<2-> \textbf{Finding 2:} Significant improvement in efficiency (2x faster)
        \item<3-> \textbf{Finding 3:} Generalizes well across different datasets
        \item<4-> \textbf{Finding 4:} Robust to noise and variations
    \end{itemize}
    
    \onslide<5->{
        \begin{block}{Impact}
            These results demonstrate the effectiveness of our approach.
        \end{block}
    }
\end{frame}

% ============================================
% SECTION 5: DISCUSSION
% ============================================

\section{Discussion}

\begin{frame}{Discussion}
    \begin{block}{Strengths}
        \begin{itemize}
            \item Strength 1: High accuracy
            \item Strength 2: Computational efficiency
            \item Strength 3: Practical applicability
        \end{itemize}
    \end{block}
    
    \vspace{0.3cm}
    
    \begin{alertblock}{Limitations}
        \begin{itemize}
            \item Limitation 1: Requires large training data
            \item Limitation 2: Limited to specific domains
        \end{itemize}
    \end{alertblock}
\end{frame}

\begin{frame}{Future Work}
    \textbf{Planned extensions:}
    \begin{enumerate}
        \item Extend to multi-modal data
        \item Investigate theoretical properties
        \item Apply to real-world scenarios
        \item Develop user-friendly tools
    \end{enumerate}
\end{frame}

% ============================================
% SECTION 6: CONCLUSION
% ============================================

\section{Conclusion}

\begin{frame}{Conclusion}
    \begin{block}{Summary}
        \begin{itemize}
            \item Addressed the problem of [X]
            \item Proposed a novel approach based on [Y]
            \item Demonstrated effectiveness through experiments
            \item Achieved state-of-the-art results
        \end{itemize}
    \end{block}
    
    \vspace{0.3cm}
    
    \begin{block}{Contributions}
        \begin{enumerate}
            \item Theoretical contribution: [X]
            \item Methodological contribution: [Y]
            \item Empirical contribution: [Z]
        \end{enumerate}
    \end{block}
\end{frame}

% ============================================
% THANK YOU SLIDE
% ============================================

\begin{frame}[plain]
    \begin{center}
        \Huge{\textbf{Thank You!}}
        
        \vspace{1cm}
        
        \Large{Questions?}
        
        \vspace{1cm}
        
        \normalsize
        Your Name \\
        \texttt{your.email@example.com} \\
        
        \vspace{0.5cm}
        
        \small
        Slides available at: \url{https://your-website.com/slides}
    \end{center}
\end{frame}

% ============================================
% BACKUP SLIDES
% ============================================

\appendix

\begin{frame}{Backup: Additional Details}
    \begin{itemize}
        \item Additional technical details
        \item Extra experimental results
        \item Supplementary information
    \end{itemize}
\end{frame}

\begin{frame}{Backup: References}
    \begin{thebibliography}{9}
        \bibitem{ref1} Author A. et al. (2023). "Paper Title". Conference/Journal.
        \bibitem{ref2} Author B. et al. (2022). "Another Paper". Conference/Journal.
        \bibitem{ref3} Author C. et al. (2021). "Related Work". Conference/Journal.
    \end{thebibliography}
\end{frame}

\end{document}
